\chapter{Applications of Simple Linear Regression in Finance}
\label{chap:Applications of Simple Linear Regression in Finance}

\begin{center}
\begin{tabular}{p{0.06\linewidth} | p{0.88\linewidth}}
	\toprule
	\multicolumn{2}{l}{\textbf{LEARNING OUTCOMES}} \\
	\midrule
	a & describe, interpret, and explain simple linear regression, including coefficient estimation using the least squares criterion \\
	\midrule
	b & describe and compare the assumptions of simple linear regression, identify violations through analyzing residuals, evaluate the estimated mode's goodness-of-fit and regression coefficients, and results of ANOVA estimates \\
	\midrule
	c & calculate and interpret predicted values, the standard error of the estimate, and prediction intervals for the dependent variable in a simple linear regression model and describe different functional forms \\
	\midrule
	d & calculate and interpret the variable estimates of the capital asset pricing model (CAPM) \\
	\bottomrule
\end{tabular}
\end{center}

\vspace{1em}

\begin{itemize}
	\item[] \textbf{Linear Regression Estimation}
	      \begin{itemize}
			\item[] Linear Regression Analysis
			\item[] Coefficient Estimation of Linear Regression
			\item[] Estimation Residuals
			\item[] Data Types
	      \end{itemize}
	\item[] \textbf{Limitations of Linear Regression Models}
	      \begin{itemize}
			\item[] Assumptions of the Simple Linear Regression Model
			\item[] Residual Analysis
			\item[] Evaluating Goodness-of-Fit and Regression Coefficients
			\item[] Hypothesis Testing
	      \end{itemize}
	\item[] \textbf{Using Linear Regression Models for Prediction}
	      \begin{itemize}
			\item[] Prediction Using a Simple Regression
			\item[] Modeling Non-Linear Relationships in Economic and Financial Data
			\item[] Practical Implementation
			\item[] Hypothesis Testing
	      \end{itemize}
	\item[] \textbf{Using Linear Regression Models to Evaluate Financial Assets: The Capital Asset Pricing Model}
\end{itemize}
